% =============================================================================
% ANNEXE L : Note de pivot du contrat d'apprentissage
% =============================================================================
% Reconstitution datée du basculement du 3e critère du contrat d'apprentissage
% (design / lecture de marché -> business layer / gouvernance de projet),
% conformément à l'exigence Tardif d'autorégulation tracée.
% =============================================================================

\chapter{Note de pivot --- glissement du 3\textsuperscript{e} critère du contrat d'apprentissage}
\label{annex:note-pivot}

\annexefiche
  {Note rédigée en mai 2026 (chronologie reconstituée a posteriori)}
  {Documentation réflexive du pivot entre contrat initial et trajectoire effectivement observée}
  {Autorégulation méthodologique : justification du glissement vers la \textit{business layer}}
  {Aucun caviardage (document réflexif ; limite assumée sur l'absence de trace contemporaine)}

\begin{tcolorbox}[colback=yellow!8, colframe=orange!70!black, boxrule=0.5pt, arc=3pt,
  left=8pt, right=8pt, top=6pt, bottom=6pt,
  title={\textbf{\small Statut de cette annexe}},
  fonttitle=\small\bfseries, coltitle=white, attach boxed title to top left={yshift=-2mm, xshift=4mm},
  boxed title style={colback=orange!70!black, arc=3pt, boxrule=0pt}]
\small
Cette note est une \textbf{reconstitution \emph{a posteriori}}, produite au moment de la rédaction finale du portfolio. Elle n'a pas été archivée en temps réel comme un point d'étape formel avec le professeur référent. Son statut est donc celui d'une \emph{rétrospective d'autorégulation} et non d'une trace contemporaine. Cette limite est explicitement assumée : elle est mentionnée dans la note méthodologique du chapitre~\ref{chap:but_portfolio} parmi les biais reconnus du dispositif, et identifiée par l'audit académique comme une faiblesse à corriger pour les cycles futurs (mise en place d'un journal de bord daté en cours d'année).
\end{tcolorbox}

\bigskip

\section*{1. Objet}

La présente note documente le basculement du \textbf{3\textsuperscript{e} critère du contrat individuel d'apprentissage} (cf.~annexe~\ref{annex:contrat-apprentissage}), initialement formulé en septembre 2025 comme \emph{« développer une conception centrée utilisateur et une meilleure lisibilité des solutions proposées »} (axe design / lecture de marché), vers un nouvel axe : \textbf{passage progressif vers la \textit{business layer}\footnote{Voir la note terminologique~\ref{fn:business-layer} dans le chapitre \textit{Perspectives}.} et la gouvernance de projet}, soutenu par la certification PRINCE2~\citep{prince2_2017}.

L'argumentation détaillée de ce pivot figure dans la section \emph{Pourquoi ce glissement par rapport au contrat d'apprentissage initial ?} du chapitre~\ref{chap:perspectives}. La présente annexe en restitue la chronologie, les déclencheurs et l'autorégulation associée.

\section*{2. Chronologie reconstituée}

\begin{tcolorbox}[colback=gray!5, colframe=gray!50, boxrule=0.4pt, arc=3pt,
  left=8pt, right=8pt, top=6pt, bottom=6pt]
\small
Les dates qui suivent sont reconstituées à partir des artefacts produits (note de décision, backlog, feuille de route, certification PRINCE2) et de la mémoire de l'auteur. Elles n'ont pas la précision d'un journal de bord tenu en temps réel.
\end{tcolorbox}

\smallskip

\begin{itemize}
  \item \textbf{Septembre 2025} --- Rédaction et validation du contrat d'apprentissage avec le professeur référent. Les trois axes retenus sont : (1) collaboration / délégation, (2) argumentation de valeur et lecture de marché, (3) conception centrée utilisateur et lisibilité des solutions. L'axe \emph{business layer / gouvernance} n'est \emph{pas} explicité comme tel ; il est latent dans l'axe~2.
  \item \textbf{Octobre--décembre 2025} --- Premiers mois de formation. Les cours \emph{Gestion de projet 1}, \emph{Méthodes de recherche 2} et \emph{Médiation scientifique} fournissent des outils que je commence à transposer dans l'environnement TWIICE. La première version de la note de décision technique (NDT-SW-001, annexe~\ref{annex:ndt-sw-001}) est produite. Première intuition : les compétences réellement mobilisées dans la pratique professionnelle relèvent davantage de la \emph{gouvernance décisionnelle} que du design d'interface.
  \item \textbf{Janvier 2026} --- Premier sprint adapté IEC~62304 chez TWIICE ; production du Software Development Plan (annexe~\ref{annex:sdp-markdown}). Constat : les artefacts qui produisent le plus de valeur dans ma pratique sont des artefacts de \emph{cadrage} (SDP, backlog priorisé, roadmap), non des artefacts de conception d'interface.
  \item \textbf{Février 2026} --- Préparation et passage de la certification PRINCE2~\citep{prince2_2017}, hors du périmètre prévu par le contrat d'apprentissage. Le passage se fait sur initiative personnelle, en lien avec un besoin perçu de \emph{cadre de gouvernance} pour structurer les arbitrages projet.
  \item \textbf{Mars 2026} --- \textbf{Point de bascule.} La conjonction des trois signaux (cohérence des artefacts produits, obtention de la certification PRINCE2, lecture en cours du chapitre Perspectives en cours de rédaction) fait apparaître que l'axe~3 du contrat d'apprentissage n'est plus le bon \emph{intitulé} de ce qui se développe réellement. Le mot juste n'est pas « design » mais \emph{gouvernance de projet et passage à la business layer}.
  \item \textbf{Avril--mai 2026} --- Documentation explicite de ce glissement dans le chapitre~\ref{chap:perspectives} du portfolio, et formalisation rétroactive de la présente note. Le contrat d'apprentissage initial n'est pas modifié (il reste l'engagement contractuel signé en septembre) ; le portfolio assume le décalage en l'explicitant.
\end{itemize}

\section*{3. Déclencheurs identifiés}

Trois éléments concourants ont rendu le pivot nécessaire :

\begin{enumerate}
  \item \textbf{Convergence des artefacts.} Les preuves accumulées au fil des compétences A2, B4 et C2 (note de décision, backlog priorisé, SDP, rétrospective, matrice de priorisation, roadmap) sont \emph{toutes} des artefacts de cadrage et de gouvernance. Aucune n'est un artefact de design d'interface. La lecture rétrospective de cette série révèle un axe que le contrat initial n'avait pas nommé.
  \item \textbf{Signal externe convergent (PRINCE2).} La préparation de PRINCE2, engagée pour des raisons indépendantes du portfolio, fait apparaître un alignement fort entre ses concepts (gouvernance par étapes, gestion des risques, tolérances, arbitrages) et la posture en construction.
  \item \textbf{Lecture critique du contrat initial.} L'axe~3 (« conception centrée utilisateur et lisibilité des solutions ») est apparu \emph{a posteriori} comme une formulation \emph{symptomatique} d'un manque plus structurant : non pas l'absence de capacité de design, mais l'absence de cadre de pilotage qui aurait permis à un éventuel design de trouver sa place dans une chaîne de décision crédible.
\end{enumerate}

\section*{4. Conséquences sur l'évaluation du contrat}

\textbf{Ce qui est conservé.} Les axes~1 (collaboration / délégation) et~2 (argumentation de valeur) du contrat sont maintenus tels quels. Les trois compétences certifiées (A2, B4, C2) restent celles définies dans le contrat. La logique d'alignement entre objectifs, compétences et choix de cours documentée au chapitre~\ref{chap:contrat_apprentissage} reste pertinente.

\textbf{Ce qui est déplacé.} L'axe~3 est \emph{recadré}, non remplacé. L'intention sous-jacente (renforcer une dimension qui me manquait à l'entrée en formation) est inchangée ; sa formulation devient plus juste à la lumière des apprentissages effectivement réalisés. Le design et la lecture de marché restent des angles de progression secondaires, non prioritaires pour la suite immédiate.

\textbf{Ce qui est explicitement assumé comme limite méthodologique.} Cette note de pivot n'est pas un point d'étape contractualisé en cours d'année. Elle ne porte pas la signature du professeur référent ; elle ne résulte pas d'une revue formelle datée. Pour les cycles d'apprentissage futurs, la mise en place d'un \textbf{journal de bord daté} et d'un \textbf{point d'étape semestriel contractualisé} constitueraient les améliorations méthodologiques permettant de produire des traces d'autorégulation \emph{contemporaines} plutôt que \emph{rétrospectives}.

\section*{5. Inscription dans le cadre Tardif}

Au sens de~\citet{tardif2006}, l'autorégulation est l'une des dimensions constitutives du portfolio d'apprentissage : l'étudiant doit pouvoir documenter \emph{comment} il ajuste sa trajectoire à la lecture de ses apprentissages effectifs. Cette annexe répond à cette exigence en restituant explicitement le glissement, ses déclencheurs et son arbitrage, plutôt que de le laisser apparaître seulement \emph{en creux} entre le contrat initial (annexe~\ref{annex:contrat-apprentissage}) et le chapitre~\ref{chap:perspectives}. La limite (reconstitution \emph{a posteriori} plutôt que trace contemporaine) est tenue comme l'un des biais explicitement assumés du dispositif, conformément à la note méthodologique du chapitre~\ref{chap:but_portfolio}.



